% easychair.tex,v 3.5 2017/03/15

\documentclass{easychair}

\usepackage{doc}

% use this if you have a long article and want to create an index
% \usepackage{makeidx}

% In order to save space or manage large tables or figures in a
% landcape-like text, you can use the rotating and pdflscape
% packages. Uncomment the desired from the below.
%
% \usepackage{rotating}
% \usepackage{pdflscape}

\makeatletter
\let\code\relax
\let\endcode\relax
\makeatother

\usepackage{agda}
\usepackage{newunicodechar}
\newunicodechar{⇒}{\ensuremath{\Rightarrow}}
\newunicodechar{◃}{\ensuremath{\triangleleft}}
\newunicodechar{▹}{\ensuremath{\triangleright}}
\newunicodechar{₁}{\ensuremath{_1}}
\newunicodechar{∈}{\ensuremath{\in}}
\newunicodechar{⊤}{\ensuremath{\top}}
\newunicodechar{≅}{\ensuremath{\cong}}
\newunicodechar{ℕ}{\ensuremath{\nat}}

%\makeindex

%% Front Matter
%%
% Regular title as in the article class.
%
\title{Containers in Higher Kinds}

% Authors are joined by \and. Their affiliations are given by \inst, which indexes
% into the list defined using \institute
%
\author{
  Thorsten Altenkirch\inst{1} \and
  H{\aa}kon Robbestad Gylterud\inst{2} \and
  Zhili Tian\inst{1}
}

% Institutes for affiliations are also joined by \and,
\institute{
  $^{1}$School of Computer Science, University of Nottingham\\
  \email{\{psztxa,psxzt8\}@nottingham.ac.uk}\\
  $^{2}$Department of Informatics, University of Bergen\\
  \email{hakon.gylterud@uib.no}
}

%  \authorrunning{} has to be set for the shorter version of the authors' names;
% otherwise a warning will be rendered in the running heads. When processed by
% EasyChair, this command is mandatory: a document without \authorrunning
% will be rejected by EasyChair

\authorrunning{Altenkirch, Gylterud, Tian}

% \titlerunning{} has to be set to either the main title or its shorter
% version for the running heads. When processed by
% EasyChair, this command is mandatory: a document without \titlerunning
% will be rejected by EasyChair
\titlerunning{Containers in Higher Kinds}

\begin{document}

\maketitle

\begin{abstract}
We generalize the initial algebra semantics from inductive types to inductive 
functors such as list, and to higher-order inductive functors such as monad 
transformers, by defining a notion of higher-order containers. We begin by 
investigating second-order case and show that every container functor can 
be uniquely represented as a second-order W-type. We then define a syntax for
higher-order containers whose semantics correspond to hereditary functors.
\end{abstract}

\begin{code}[hide]
{-# OPTIONS --guardedness --allow-unsolved-metas #-}
module main where
\end{code}

\section{Introduction}

\textbf{Inductive types} are fundamental constructs in functrional programming 
and in type theory, allowing the definition of boolean, natural numbers, lists,
trees, etc. through a finite set of constructors. An inductive type is 
\textbf{strictly positive} if it appears only on positive positions in its 
constructors. Focusing solely on strictly positive inductive types, we note that 
all such can be described as the initial algebra of \textbf{containers}.

\subsection{Containers}

A container is given by a pair of shape \texttt{S : Set} and positions 
\texttt{P : S → Set}.

TODO: 
1. Give the definition and the intuition of containers.
2. Give example and general W-type (and M-type); insert diagrams.
3. Talk about indexed inductive types and indexed containers; example is Fin. 
4. What about types that are functors? Traditionally we fix A : Set and say
List A is W-type. But we can do better.

\subsection{Setting}

TODO:
1. It is formalized in Agda and we have funExt, setExt, uip.

make a random citation\cite{bird1998nested}

\section{Second-order Containers}

TODO:
1. Give the definition of 2Cont
2. Structures and closure properties of 2Cont.
3. Semantics: (Set → Set) → Set → Set, Func ⇒ Func, Cont ⇒ Cont.
4. 2W : 2Cont → Func, fold2W : ...; give code and diagram.
5. Talk about 2M.
6. Is W still enough?

\section*{Higher-order Containers}

\bibliographystyle{plain}
\bibliography{references}

%------------------------------------------------------------------------------
\end{document}